\documentclass[10pt,a4paper]{amsart}
\usepackage[margin=19mm]{geometry}
\usepackage{amsmath,amssymb,mathtools}
\usepackage[scr=rsfs,cal=euler]{mathalfa}
\usepackage{xfrac}
\usepackage[unicode,hidelinks,bookmarks=false]{hyperref}
\newcommand{\ie}{i.e.\ }
\newcommand{\cf}{cf.\ }
\newcommand{\resp}{respectively\ }
\newcommand{\Subsection}{\subsection}
\providecommand{\nobreakdash}{\nobreak\hbox{-}}
\setlength{\emergencystretch}{2em}
\multlinegap0pt
\usepackage{bm}
\usepackage{bbm}
\usepackage{dsfont}
\usepackage{xspace}
\usepackage{cancel}
\usepackage{mathtools}
\usepackage{amsthm}
\makeatletter
\@ifundefined{theorem}{\newtheorem{theorem}{Theorem}}{}
\@ifundefined{lemma}{\newtheorem{lemma}[theorem]{Lemma}}{}
\@ifundefined{proposition}{\newtheorem{proposition}[theorem]{Proposition}}{}
\makeatother
\DeclareMathOperator{\im}{im}
\DeclarePairedDelimiter{\abs}{\lvert}{\rvert}
\DeclarePairedDelimiter{\parens}{\lparen}{\rparen}
\DeclarePairedDelimiter{\set}{\{}{\}}
\newcommand\Q{{\mathbb Q}}
\newcommand\Z{{\mathbb Z}}
\title[Hilbert Series for Configuration Spaces of Punctured Surface]{Hilbert Series for Configuration Spaces of Punctured Surfaces}
\author{Yifeng Huang, Eric Ramos}
\date{Source: arXiv:2507.09746v1 (CC BY 4.0)}
\begin{document}
\maketitle

\noindent\emph{Source and licence.} Hilbert Series for Configuration Spaces of Punctured Surfaces. Version arXiv:2507.09746v1 (CC BY 4.0), distributed under the Creative Commons Attribution 4.0 International licence, as recorded in the arXiv licence metadata for this exact version and shown on the abstract page https://arxiv.org/abs/2507.09746v1. Full rights notice: licenses/arxiv-2507.09746-CC-BY-4.0.txt.\par\smallskip

\noindent\textit{Editorial scope.} Counted unit: the Proposition whose source label is \texttt{eq:hilbert-hgn} in the article (source0.tex lines 700--712, source label \texttt{eq:hilbert-hgn}), delivered together with the article's own complete original proof of it (source0.tex lines 714--759, one \texttt{proof} environment of 46 lines). No mathematics is written, repaired or re-derived: statement and proof are the article's own text line for line. Required context retained uncounted: eq:grading (lines 522--524); eq:TheC's (lines 544--546); lem:hard-lefschetz (lines 686--689). Every definition, display and label of those lines is retained unchanged. Omitted from this package and not redistributed: 1--521, 525--543, 547--685, 690--699, 713--713, 760--985. Nothing inside a retained excerpt is omitted or altered: every retained body is delivered exactly as the article's lines are written, so no omission receipt of any kind accompanies this package. Citations: the retained excerpts carry no citation command at all, so no bibliography entry of the article is transcribed and no reference is redistributed. Reference adaptation: none was needed and none was made; every reference inside the delivered unit resolves inside it, so no unresolved reference or citation is produced. Printed slips kept literal: no printed word, symbol, hypothesis or number of the counted statement or of its proof was altered. Layout and class: the article's own class is \texttt{amsproc} with packages including geometry, bbm, amsfonts, amsmath, amsthm, amssymb, latexsym, mathrsfs, stmaryrd, hyperref, mathtools, graphicx, subcaption, tikz; the wrapper is the lane's pinned \texttt{amsart} editorial wrapper at 10pt on A4, which restates the theorem-like environments the retained text needs and adds the article's own macro definitions that the retained text needs (\texttt{\textbackslash Q}, \texttt{\textbackslash Z}, \texttt{\textbackslash abs}, \texttt{\textbackslash im}, \texttt{\textbackslash parens}, \texttt{\textbackslash set}), with extra wrapper packages (bm, bbm, dsfont, xspace, cancel, mathtools). The delivered numbering is the unit's own. Assets: no retained span contains a figure environment, a graphics-inclusion command, a PDF-page inclusion, a table or a TikZ picture; no image file is copied into this package, the package declares no asset dependency and no image file is redistributed with it. Licence: version arXiv:2507.09746v1 of this article is distributed under the Creative Commons Attribution 4.0 International licence, as recorded in the arXiv licence metadata for this exact version and shown on the abstract page https://arxiv.org/abs/2507.09746v1. A full rights notice is delivered at licenses/arxiv-2507.09746-CC-BY-4.0.txt. Characters: every retained excerpt is pure ASCII, so the delivered engine, XeLaTeX with its default Latin Modern text font, reports no missing character.
\begin{equation}\label{eq:grading}
    p(\mathbf{t})=\abs{S}+\abs{T}+\abs{\vec u}+\abs{\vec v}, \; q(\mathbf{t})=r+\abs{\vec u}+\abs{\vec v}, \; w_1(\mathbf{t})=r+\abs{S}+2\abs{\vec u}+\abs{\vec v}, \; w_2(\mathbf{t})=r+\abs{T}+\abs{\vec u}+2\abs{\vec v}.
\end{equation}

\begin{equation} \label{eq:TheC's}
    -\frac12 df_{(1,S,T,\vec u,\vec v)}=\sum_{\substack{S',T'\\\abs{S'}+\abs{T'}=\abs{S}+\abs{T}+2}} c^{S,T}_{S',T'}f_{(1,S',T',\vec u,\vec v)}.
\end{equation}

\begin{lemma}
    The operator $\iota:\Omega^i\to \Omega^{i+2}$ is injective if $i\leq g-1$ and surjective if $i\geq g-1$. 
    \label{lem:hard-lefschetz}
\end{lemma}

\begin{proposition}
    We have
    \begin{equation}\label{eq:hilbert-hgn}
        \sum_{p,q,w_1,w_2,n} \dim H_{g,n}^{(p,q,w_1,w_2)} x^{w_1}y^{w_2}(-u)^{p+q}t^n =\frac{\Phi_g\set{(1-xyz^2)(1-xz)^g(1-yz)^g}}{(1-t)(1-x^2 yu^2t^2)^g(1-xy^2u^2t^2)^g},
    \end{equation}
    where $\Phi_g$ is a $\Q[x,y]$-linear map defined by 
    \begin{equation}\label{eq:shiftDef}
    \Phi_g(z^j)=\begin{cases}
    u^j t^j,& 0\leq j\leq g;\\
    u^{j-1}t^j,& g+2\leq j\leq 2g+2.
    \end{cases}
    \end{equation}
\end{proposition}

\begin{proof}

Given $\rho\in \set{0,1}$, $0\leq \sigma,\tau\leq g$, $\vec u,\vec v\in \Z_{\geq 0}^g$, define $V_{(\rho,\sigma,\tau,\vec u,\vec v)}=\mathrm{Span}_\Q\set{f_{(\rho,S,T,\vec u,\vec v)}: \abs{S}=\sigma, \abs{T}=\tau}$. Consider
\begin{equation}
    H_{(1,\sigma,\tau,\vec u,\vec v)}=\ker(d:V_{(1,\sigma,\tau,\vec u,\vec v)}\to V_{(0,\sigma+1,\tau+1,\vec u,\vec v)})
\end{equation}
and
\begin{equation}
    H_{(0,\sigma,\tau,\vec u,\vec v)}=\frac{V_{(0,\sigma,\tau,\vec u,\vec v)}}{\im(d:V_{1,\sigma-1,\tau-1,\vec u,\vec v}\to V_{0,\sigma,\tau,\vec u,\vec v})}.
\end{equation}

Then
\begin{equation}
    H_{g,n}^{(p,q,w_1,w_2)}=\bigoplus_{\rho,\sigma,\tau,\vec u,\vec v}H_{(\rho,\sigma,\tau,\vec u,\vec v)}
\end{equation}
where the sum ranges over all $(\rho,\sigma,\tau,\vec u,\vec v)$ such that \eqref{eq:grading} holds and $2\rho+\sigma+\tau+\abs{\vec u}+\abs{\vec v}\leq n$. 

By Lemma \ref{lem:hard-lefschetz} and the fact that the same structure constants $c^{S,T}_{S',T'}$ appear in $\iota$ and \eqref{eq:TheC's}, we have
\begin{equation}
    \dim H_{(1,\sigma,\tau,\vec u,\vec v)}=
    \begin{cases}
        0, & \sigma+\tau\leq g-1,\\
        \binom{g}{\sigma}\binom{g}{\tau}-\binom{g}{\sigma+1}\binom{g}{\tau+1}, & \sigma+\tau\geq g-1,
    \end{cases}
\end{equation}
and
\begin{equation}
    \dim H_{(0,\sigma,\tau,\vec u,\vec v)}=
    \begin{cases}
        \binom{g}{\sigma}\binom{g}{\tau}-\binom{g}{\sigma-1}\binom{g}{\tau-1}, & \sigma+\tau\leq g+1\\
        0, &\sigma+\tau\geq g+1.
    \end{cases}
\end{equation}

Summing up the contribution to the Hilbert series, and noting that each $(\rho,\sigma,\tau,\vec u, \vec v)$ contributes to $H_{g,n}$ if and only if $2\rho+\sigma+\tau+\abs{\vec u}+\abs{\vec v}\leq n$, we get
\begin{align}
    &\sum_{p,q,w_1,w_2,n} \dim H_{g,n}^{(p,q,w_1,w_2)} x^{w_1}y^{w_2}(-u)^{p+q}t^n\\ 
    &= \sum_{\sigma,\tau,\vec u,\vec v} \dim H_{(1,\sigma,\tau,\abs{\vec u},\abs{\vec v)}} x^{1+\sigma+2\abs{\vec u}+\abs{\vec v}} y^{1+\tau+\abs{\vec u}+2\abs{\vec v}}(-u)^{1+\sigma+\tau+2\abs{\vec u}+2\abs{\vec v}}t^{2+\sigma+\tau+2\abs{\vec u}+2\abs{\vec v}}/(1-t)+\\
    &\sum_{\sigma,\tau,\vec u,\vec v} \dim H_{(0,\sigma,\tau,\vec u,\vec v)} x^{\sigma+2\abs{\vec u}+\abs{\vec v}} y^{\tau+\abs{\vec u}+2\abs{\vec v}}(-u)^{\sigma+\tau+2\abs{\vec u}+2\abs{\vec v}}t^{\sigma+\tau+2\abs{\vec u}+2\abs{\vec v}}/(1-t)\\
    &=\frac{1}{(1-t)(1-x^2 yu^2t^2)^g(1-xy^2u^2t^2)^g}\parens*{-xyut^2 P(-xut,-yut)+Q(-xut,-yut)},
\end{align}
where $$P(x,y):=\sum_{\substack{\sigma,\tau\geq 0\\\sigma+\tau\geq g-1}} x^\sigma y^\tau \parens*{\binom{g}{\sigma}\binom{g}{\tau}-\binom{g}{\sigma+1}\binom{g}{\tau+1}}$$ and $$Q(x,y):=\sum_{\substack{\sigma,\tau\geq 0\\\sigma+\tau\leq g}}x^\sigma y^\tau \parens*{\binom{g}{\sigma}\binom{g}{\tau}-\binom{g}{\sigma-1}\binom{g}{\tau-1}}$$
are the generating functions for $\dim H_{(1,\sigma,\tau,\abs{\vec u},\abs{\vec v})}$ and $\dim H_{(0,\sigma,\tau,\vec u,\vec v)}$, respectively.

The proof is complete once we prove the following lemma.
\end{proof}

\end{document}
